\emph{The set of values.} By clause (ii) of Theorem \claimref{synthesis.chsh-over-the-rationals-pure-states-on-a-circle-and-the-pell-records} the pure states of the sector have $\omega(C) = X$, $\omega(C^{\perp}) = Y$ with $(X, Y)$ a rational point of $X^{2} + Y^{2} = 8$, and every rational point occurs; so the set of their values of $C$ is $\mathcal{V}$. It is symmetric under $X \mapsto -X$ and does not contain $0$ ($8$ is not a rational square), so it is enough to treat the positive values.

\emph{(i).} Let $X = P/Q \in \mathcal{V}$ in lowest terms, $P, Q \ge 1$. Then $8Q^{2} - P^{2}$ is the square of a rational and an integer, so $8Q^{2} - P^{2} = R^{2}$ with $R \ge 1$ an integer ($R = 0$ would make $8$ a square). A sum of two squares divisible by $4$ has both terms even, so $P = 2a$, $R = 2b$ and $a^{2} + b^{2} = 2Q^{2}$. Since $P$ is even and the fraction is in lowest terms, $Q$ is odd; then $a^{2} + b^{2} \equiv 2$ modulo $4$, and $a$, $b$ are odd. $\gcd(a, Q) = 1$ since $\gcd(P, Q) = 1$. Let $\ell$ be a prime dividing $b$. Then $a^{2} \equiv 2Q^{2}$ modulo $\ell$, and $\ell \nmid Q$ (otherwise $\ell \mid a$), so $2$ is a square modulo $\ell$. For $\ell = 3$ and $\ell = 5$ it is not (the squares are $\{0, 1\}$ and $\{0, 1, 4\}$), so $b$ is odd and divisible by neither $3$ nor $5$: $b = 1$ or $b \ge 7$. (By the supplementary law every prime factor of $b$ is $\pm 1$ modulo $8$; only the cases $3$ and $5$ are used.) The deficit is $8 - X^{2} = R^{2}/Q^{2} = 4b^{2}/Q^{2}$. $b = 1$ is $a^{2} - 2Q^{2} = -1$, the equation of clause (v) of the stored theorem, whose positive solutions are the $(x_{n}, Q_{n})$; so $b = 1$ exactly for $X = S_{n}$.

\emph{(ii).} $x_{n}^{2} = 2Q_{n}^{2} - 1 < 2Q_{n}^{2}$, so $Q_{n+1} = 2x_{n} + 3Q_{n}$ satisfies $Q_{n+1}^{2} = 4x_{n}^{2} + 12x_{n}Q_{n} + 9Q_{n}^{2} < 17Q_{n}^{2} + 12x_{n}Q_{n}$, and $(12x_{n}Q_{n})^{2} = 144x_{n}^{2}Q_{n}^{2} < 288Q_{n}^{4} < (17Q_{n}^{2})^{2}$; hence $Q_{n+1}^{2} < 34Q_{n}^{2}$. Let $X = P/Q \in \mathcal{V}$ be positive, in lowest terms, with $b \ge 7$. Then $2Q^{2} = a^{2} + b^{2} \ge 50$, so $Q > 1 = Q_{0}$; let $n$ be the largest index with $Q_{n} < Q$. Then $Q \le Q_{n+1}$, so $Q^{2} < 34Q_{n}^{2}$ and
\[
8 - S_{n}^{2} = \frac{4}{Q_{n}^{2}} < \frac{136}{Q^{2}} < \frac{196}{Q^{2}} \le \frac{4b^{2}}{Q^{2}} = 8 - X^{2},
\]
so $S_{n} > X$ (both are positive).

\emph{(iii).} The $S_{n}$ increase with $n$ (their deficits $4/Q_{n}^{2}$ decrease). Fix $D \ge 1$ and let $n$ be the largest index with $Q_{n} \le D$. A positive value of denominator at most $D$ is either some $S_{m}$ with $Q_{m} \le D$, hence $m \le n$ and $S_{m} \le S_{n}$, or has $b \ge 7$ and by (ii) is smaller than some $S_{m}$ with $Q_{m}$ smaller than its denominator, hence again smaller than $S_{n}$. So the largest value of denominator at most $D$ is $S_{n}$, attained only there. A value is a running maximum when it is the largest value of denominator at most its own denominator; by the above these are exactly the $S_{n}$.
